\par\noindent\textbf{(1)} 守恒能量 m l²θ̇²/2+mgl(1−cosθ)=mgl(1−cosθ₀)，因此 θ̇²=2ω₀²(cosθ−cosθ₀)。四分之一周期 T/4=∫₀\textasciicircum{}\{θ₀\}dθ/[√2ω₀√(cosθ−cosθ₀)]，非线性振动速度在端点为零但该积分对θ₀<π可收敛。
\par\noindent\textbf{(2)} 令 k=sin(θ₀/2)，再以 sin(θ/2)=k sinφ，注意 cosθ−cosθ₀=2k²cos²φ、dθ=2kcosφ dφ/√(1−k²sin²φ)，得到 T=4K(k)/ω₀，K(k)=∫₀\textasciicircum{}\{π/2\}(1−k²sin²φ)\textasciicircum{}\{−1/2\}dφ。
\par\noindent\textbf{(3)} 使用二项式展开，K(k)=(π/2)[1+k²/4+9k⁴/64+O(k⁶)]；又 k²=θ₀²/4−θ₀⁴/48+…，故 T=(2π/ω₀)[1+θ₀²/16+11θ₀⁴/3072+O(θ₀⁶)]。大振幅使周期变长，小角极限为线性摆周期。
\par\noindent\textbf{(4)} 当 θ₀→π⁻，k→1⁻，椭圆积分 K(k)∼ln[4/√(1−k²)]；于是T以对数发散。靠近顶部的线性化θ̈≈+ω₀²(π−θ)对应指数离开不稳定平衡，粒子在分离轨道附近停留极久；精确θ₀=π且初速零时永久静止而不构成周期振动。
